\chapter{智慧水利概述与平台架构基础}

\paragraph*{学习目标}

学完本章，读者应能够：

\begin{enumerate}
\item 说明智慧水利的业务目标、能力边界及其与一般信息化系统的区别；
\item 解释感知、平台、应用和用户四层之间的数据上行与控制下行关系；
\item 区分测站通信规约、消息传输协议和数据交换标准的用途；
\item 使用软件工程生命周期理解智慧水利平台的建设与持续运行；
\item 沿一次洪水预警链路识别数据、模型、人员和审计证据。
\end{enumerate}

\paragraph*{引言}

智慧水利平台不是把地图、图表和三维模型放在同一个页面，也不是用某一种新技术替换原有系统。它以水利业务目标为牵引，把监测感知、通信、数据治理、模型计算、业务协同和运行管理连接为可追溯闭环。本章先建立全书共同语言，再用一次洪水预警说明平台如何工作。

\section{智慧水利的概念与发展}

\subsection{定义、边界与基本特征}

本教材把智慧水利理解为：以水利对象和治理活动为服务对象，以可信数据为基础，以模型和知识为支撑，通过数字化场景、智慧化模拟和精准化决策提升预报、预警、预演、预案能力的业务与技术体系。它包含软件平台，也包含测站、网络、制度、人员和工程运行边界。

智慧水利与“建设一个信息系统”的差异主要体现在五个方面：

\begin{table}[htbp]
\centering
\caption{智慧水利平台的基本特征}
\label{tab:smart-water-features}
\begin{tabular}{p{2.4cm}p{5.0cm}p{5.2cm}}
\toprule
特征 & 含义 & 可核验表现 \\
\midrule
业务牵引 & 从防洪、供水、工程安全等任务反推能力 & 每项功能对应角色、场景和验收问题 \\
数据可信 & 统一对象编码、时间、单位、质量码和来源 & 任一指标可追溯到测点与处理过程 \\
模型支撑 & 模型有版本、适用范围和误差说明 & 结果能复现，越界或失败时明确提示 \\
协同闭环 & 预警进入确认、处置、复核和归档 & 责任、时限、回执和审计记录完整 \\
安全运行 & 数据、控制和平台权限按风险分级 & 高风险动作受审批约束且可回滚 \\
\bottomrule
\end{tabular}
\end{table}

工程成效必须用正确概念表达。例如，若批准的建设目标是把防洪标准由30年一遇提高到100年一遇，应表述为“防洪标准提高”；不能写成“洪灾风险从30年一遇降至100年一遇”。风险还受暴露度、脆弱性、调度和极端事件影响，不能仅由重现期替代。

\subsection{政策主线与建设方向}

2019年前后的“补短板、强监管”是水利改革发展的阶段性总基调，不能与后续智慧水利文件混写为同一政策。进入“十四五”时期，水利部把智慧水利建设纳入推动新阶段水利高质量发展的实施路径，提出“需求牵引、应用至上、数字赋能、提升能力”，并以数字孪生流域为核心建设“四预”能力\cite{smartWaterGuidance2021}。

《“十四五”期间推进智慧水利建设实施方案》提出建设七大江河数字孪生流域\cite{smartWaterGuidance2021}。随后，国家水网建设与数字孪生水利工程相关文件进一步把流域、水网和工程联系起来\cite{nationalWaterNetwork2023,twinProject2024}。本教材据此采用“业务需求—数据底板—模型平台—应用闭环”的主线，不用无出处的项目数量或成效百分比装饰结论。

政策目标不等于项目验收指标。项目仍需把“四预”分解为可测试问题，例如：输入场景是否完整、模型是否可复现、预警是否按权限流转、预演是否保留方案版本、预案是否能关联责任与资源。

\subsection{感知、通信与数据交换}

监测链路通常包含传感器、遥测终端或边缘网关、通信网络、接入服务和数据平台。不同层使用的规范不可混为一谈：

\begin{table}[htbp]
\centering
\caption{监测链路中三类规范的不同职责}
\label{tab:monitoring-protocol-roles}
\begin{tabular}{p{2.4cm}p{4.5cm}p{5.7cm}}
\toprule
类别 & 例子 & 用途 \\
\midrule
水利通信规约 & SL/T 651-2014 & 规定智能传感器与遥测终端、测站与中心站之间的数据通信\cite{slt6512014} \\
通用消息协议 & MQTT、HTTP & 在满足安全与可靠性约束时承载设备或网关消息 \\
数据交换标准 & WaterML 2.0 & 描述和交换水文观测时间序列，不等同于测站上报通信协议 \\
\bottomrule
\end{tabular}
\end{table}

技术选型取决于现场网络、功耗、时延、安全域和运维责任。涉及成本变化和性能提升时必须给出可核验基线；在线服务也不能直接证明具备离线部署能力。需要本地推理时，应先验证模型许可、算力、更新机制、数据隔离和安全评测，再决定是否采用开源权重或专用模型。

\subsection{贯穿场景：一次洪水预警如何形成闭环}

清源水库是全书使用的虚构教学工程。平台接入雨量、水位和闸门状态；本章只讨论链路与职责，不把演示参数当作真实工程结论。

汛期某时刻，库区雨量站上报分钟雨量。遥测终端按通信规约编码，接入服务校验站码、时间和单位；质量服务标记缺测、突变或回补；洪水预报模型读取同一输入快照并输出过程线；规则服务结合批准阈值生成黄色预警。值班员确认后创建工单，专业人员复核模型与现场情况，审批人决定是否启动预案，处置结果再与后续实况比较。

\begin{figure}[htbp]
\centering
\begin{tikzpicture}[node distance=6mm and 6mm,
  casebox/.style={draw=Primary,rounded corners,fill=PrimaryLight,minimum width=2.15cm,
    minimum height=.8cm,align=center,font=\small}]
\node[casebox] (sense) {雨量/水位\\监测};
\node[casebox,right=of sense] (ingest) {接入校验\\站码/时间/单位};
\node[casebox,right=of ingest] (quality) {质量标记\\原值保留};
\node[casebox,right=of quality] (model) {洪水预报\\版本/误差};
\node[casebox,right=of model] (warning) {黄色预警\\事件ID};
\node[casebox,below=of warning] (work) {确认与工单};
\node[casebox,left=of work] (review) {会商与审批};
\node[casebox,left=of review] (action) {处置与回执};
\node[casebox,left=of action] (learn) {实况复核\\模型改进};
\draw[->,draw=Primary,thick] (sense)--(ingest);
\draw[->,draw=Primary,thick] (ingest)--(quality);
\draw[->,draw=Primary,thick] (quality)--(model);
\draw[->,draw=Primary,thick] (model)--(warning);
\draw[->,draw=Primary,thick] (warning)--(work);
\draw[->,draw=Primary,thick] (work)--(review);
\draw[->,draw=Primary,thick] (review)--(action);
\draw[->,draw=Primary,thick] (action)--(learn);
\draw[->,draw=Accent,dashed,thick,bend right=28] (learn) to node[below,font=\scriptsize]{规则/模型校正} (quality);
\end{tikzpicture}
\caption{一次洪水预警从监测到复核的闭环}
\label{fig:smart-water-evolution}
\end{figure}

这个案例揭示了四条边界。第一，传感器读数经过质检后才成为可用数据；第二，模型输出是带适用范围的证据，不是自动决策；第三，平台建议不能绕过人工审批和工程规程；第四，处置完成不等于流程结束，还要用实况复核并归档。

\section{智慧水利平台架构体系}

\subsection{四层架构与双向链路}

为便于学习，本书把平台划分为感知层、平台层、应用层和用户层。分层是职责边界，不代表必须部署为四套系统。

\begin{figure}[htbp]
\centering
\begin{tikzpicture}[node distance=5mm,
  layer/.style={draw=Primary,rounded corners,fill=PrimaryLight,minimum width=11.2cm,
    minimum height=.85cm,align=center}]
\node[layer] (user) {用户层：值班员、专业分析员、审批人、运维人员};
\node[layer,below=of user] (app) {应用层：防洪预警、水资源调度、工程安全、会商与工单};
\node[layer,below=of app] (platform) {平台层：数据治理、模型服务、事件总线、空间底板、权限与审计};
\node[layer,below=of platform] (sense) {感知层：雨量、水位、流量、工情、视频、遥测终端与边缘网关};
\draw[->,draw=Primary,thick,xshift=-18mm] (sense.north)--node[left,font=\scriptsize]{观测与回执上行} (platform.south);
\draw[->,draw=Primary,thick,xshift=-18mm] (platform.north)--(app.south);
\draw[->,draw=Primary,thick,xshift=-18mm] (app.north)--(user.south);
\draw[->,draw=Accent,thick,xshift=18mm] (user.south)--node[right,font=\scriptsize]{授权指令下行} (app.north);
\draw[->,draw=Accent,thick,xshift=18mm] (app.south)--(platform.north);
\draw[->,draw=Accent,thick,xshift=18mm] (platform.south)--(sense.north);
\end{tikzpicture}
\caption{智慧水利平台四层架构及数据、控制双向链路}
\label{fig:smart-water-architecture}
\end{figure}

观测数据与设备回执自下而上进入业务应用；经授权的调度或控制意图自上而下传递。下行链路必须比一般查询具备更严格的身份、审批、有效期、设备状态校验和回滚条件。三维动画只能模拟闸门启闭与水流变化，不能替代真实控制反馈。

\subsection{架构设计原则}

\textbf{业务与数据一致。} 对象编码、单位、时间基准和质量码贯通各层；页面中的指标能够追溯到原始数据和处理规则。

\textbf{模块边界清楚。} 接入、质量、存储、模型、预警和工单各自承担稳定职责，通过版本化接口协作。某个模型或三维引擎替换时，不应重写全部业务。

\textbf{高风险操作受控。} 查询、模拟、建议、审批和执行是不同权限。平台不能因“智能”而越过工程规程和责任制度。

\textbf{失败可见并可降级。} 网络中断时缓存并标记数据龄期；模型失败时保留实时监测；三维失败时退化为二维地图、图表和表格。

\textbf{全过程可观测。} 日志、指标、追踪和审计共同回答系统是否可用、数据是否可信、模型是否运行、业务是否闭环。

这些原则在1.2节只建立总览，具体架构设计方法见第3章，前后端与消息协作见第4、5章，三维和观测数据展示见第6、7章。

\section{软件工程方法概述}

\subsection{软件及其工程特征}

软件不仅是源代码，还包括配置、数据结构、接口、脚本和运行文档。与硬件相比，软件具有以下常见工程特征：

\begin{enumerate}
\item 软件是逻辑产品，质量需要通过行为、结构和证据间接验证；
\item 开发成本主要来自分析、设计、实现、验证和协作，而非批量复制；
\item 软件不会物理磨损，但会因需求、依赖、数据和环境变化而退化；
\item 修改容易开始，却可能沿接口、数据和部署关系扩散影响；
\item 复杂度来自大量状态、组合和异常路径，不能穷尽所有输入；
\item 软件长期处于演化中，需要版本、测试、配置和运维共同治理。
\end{enumerate}

这些特征说明，智慧水利平台不能“开发完成后一次性交付”。汛期规则、测点、工程状态、模型和安全要求持续变化，团队必须保留需求依据、架构决策、测试结果和运行记录。

\subsection{生命周期与核心活动}

软件工程用系统化、可验证的方法组织软件生命周期。核心活动包括需求获取与分析、架构和详细设计、实现与集成、验证与确认、部署运行、维护演化以及最终退役\cite{ISO12207,sommerville2015}。

\begin{figure}[htbp]
\centering
\begin{tikzpicture}[node distance=6mm and 5mm,
  life/.style={draw=Primary,rounded corners,fill=PrimaryLight,minimum width=2.0cm,
    minimum height=.75cm,align=center,font=\small}]
\node[life] (req) {需求与范围};
\node[life,right=of req] (design) {架构与设计};
\node[life,right=of design] (build) {实现与集成};
\node[life,right=of build] (verify) {验证与确认};
\node[life,right=of verify] (operate) {部署与运行};
\node[life,below=of build] (evolve) {反馈、维护与演化};
\draw[->,draw=Primary,thick] (req)--(design);
\draw[->,draw=Primary,thick] (design)--(build);
\draw[->,draw=Primary,thick] (build)--(verify);
\draw[->,draw=Primary,thick] (verify)--(operate);
\draw[->,draw=Accent,thick] (operate)--(evolve);
\draw[->,draw=Accent,dashed,thick] (evolve)--(req);
\end{tikzpicture}
\caption{软件生命周期核心活动及反馈}
\label{fig:sdlc-flow}
\end{figure}

不同项目可以采用瀑布、迭代、增量、原型或敏捷方法组织活动，但不能省略责任、验证和配置控制。第2章将进一步讨论生命周期模型选型与需求分析。

\subsection{软件工程对智慧水利开发的作用}

在清源水库案例中，软件工程把模糊愿望转为可验收交付物：需求阶段明确“谁在什么场景下解决什么问题”；设计阶段确定对象、数据流和权限边界；实现阶段按接口契约协作；验证阶段覆盖正常、异常和降级路径；运行阶段用监控、工单和复盘持续改进。

软件工程也帮助团队识别“技术可行”与“工程可用”的差异。模型能运行并不代表结果可用于业务，页面能显示并不代表数据可信，设备能接收命令并不代表允许自动控制。只有将质量属性、人在回路和审计证据纳入生命周期，平台才能长期稳定运行。

\section*{小结}

智慧水利以业务目标和可信数据为基础，通过模型、应用与人员协同形成闭环。四层架构既有观测和回执上行，也有受控指令下行；通信规约、消息协议和数据交换标准承担不同职责。软件工程为需求、设计、实现、验证、运行和演化提供共同方法。

\section*{核心术语表}

\begin{longtable}{p{2.6cm}p{10.0cm}}
\toprule
术语 & 本书含义 \\
\midrule
智慧水利 & 以数字化、网络化、智能化手段提升水利治理和工程运行能力的业务与技术体系。\\
四预 & 预报、预警、预演、预案，详见第8章数字孪生案例。\\
数字孪生 & 物理对象、数字模型、数据和服务持续关联的系统方法，详见第6章6.4节。\\
三维场景 & 在统一坐标与对象编码下表达工程和环境的交互空间，详见第6章。\\
感知层 & 获取雨量、水位、流量、工情等观测及设备回执的设备与边缘资源。\\
平台层 & 提供数据、模型、事件、空间、权限和审计等共享能力的层次。\\
应用层 & 面向防洪、供水、工程安全等任务组织业务流程与界面的层次。\\
质量码 & 表示观测数据有效、缺测、异常、修正等状态的受控代码。\\
人在回路 & 把专业判断、审批和复核嵌入模型与自动化流程的机制。\\
\bottomrule
\end{longtable}

\section*{章末交付物}

完成本章后，学习小组应提交：一张带数据上行和控制下行的四层架构图；一份洪水预警链路说明；一张监测通信规约、消息协议和数据交换标准的对照表。

\section*{思考题与练习题}

\begin{enumerate}
\item 为什么智慧水利平台不能等同于三维可视化大屏？
\item 说明“防洪标准提高”与“洪灾风险降低”的区别。
\item 以图\ref{fig:smart-water-architecture}为依据，分别举出一条数据上行和控制下行链路。
\item 为什么WaterML 2.0不能直接表述为测站上报通信协议？
\item 清源水库有28个测点，每分钟每点上报1条200字节记录。忽略协议开销，计算每日原始数据量，并说明工程估算还应加入哪些开销。
\item 在“雨量上报—洪水预报—黄色预警—工单处置”链路中，列出至少四项必须保存的审计证据。
\item 选择网络中断、模型失败或三维加载失败之一，设计一个降级方案。
\item 为一个闸门启闭模拟页面说明“模拟建议”与“真实控制”之间必须设置的边界。
\end{enumerate}
